\section{Notation for the Master Chain}
\label{app:master-chain-notation}

\begin{definition}[Stage symbols of the master chain]
We use the following symbols for the principal stages of the Companion:
\begin{align*}
\mathfrak T &:= \text{triolic source},\\
\mathfrak C &:= \text{carrier layer},\\
\mathfrak C^{\mathrm{adm}} &:= \text{admissibly transported carrier},\\
\mathfrak C^{\mathrm{cl}} &:= \text{closure-stable compensated carrier},\\
\mathfrak G &:= \text{group-layer symmetry readout},\\
\mathfrak L &:= \text{null-axis structure},\\
\mathfrak P &:= \text{photon fix-sector},\\
\mathfrak R &:= \text{residual-localization sector},\\
\mathfrak H &:= \text{Higgs-mediation sector},\\
\mathfrak F &:= \text{field-equation readout}.
\end{align*}
We further write $\mathcal M_{16}$ for the Millennium matrix and
$\mathcal H_{16}$ for its global Heisenberg-compensator reading.
\end{definition}

\begin{definition}[Master chain]
The full master chain is written as
\[
\mathfrak T
\;\longrightarrow\;
\mathfrak C
\;\xrightarrow{\;\mathcal M_{16}\;}
\mathfrak C^{\mathrm{adm}}
\;\xrightarrow{\;\mathcal H_{16}\;}
\mathfrak C^{\mathrm{cl}}
\;\longrightarrow\;
\mathfrak G
\;\longrightarrow\;
\mathfrak L
\;\longrightarrow\;
\mathfrak P
\;\longrightarrow\;
\mathfrak R
\;\longrightarrow\;
\mathfrak H
\;\longrightarrow\;
\mathfrak F.
\]
When intermediate carrier detail is not needed, we also use the short form
\[
\mathfrak T
\to
\mathfrak C
\to
\mathcal M_{16}
\to
\mathcal H_{16}
\to
\mathfrak G
\to
(\mathfrak P,\mathfrak R,J\!\cdot\!\mathfrak R)
\to
\mathfrak H
\to
\mathfrak F.
\]
\end{definition}

\begin{remark}[Why this notation is useful]
This notation is useful because it prevents long later sections from drifting
into disconnected local language. Instead, one may always ask at which stage
of the master chain a given statement belongs.
\end{remark}

